\section{Related Work}
\label{sec:related}

\paragraph{Mechanistic and logic-based models of regulation.} Continuous logic models turn regulatory interactions into ordinary differential equations with Hill-type transfer functions~\cite{wittmann2009}, and cell-type-specific versions with a modified Hill transfer (TF4) and an edge-strength interpretation have been fitted to signalling data~\cite{eduati2017}. These models are interpretable but have been applied to networks of tens to hundreds of nodes; iPerturb uses the same transfer function on genome-scale networks, shares its shape parameter per regulator to keep the parameter count manageable, and makes the propagation a matrix product.

\paragraph{Biologically informed neural networks.} Genetic Neural Networks arrange one layer per gene so that each gene is predicted from its transcription-factor regulators with Hill-kinetics activations, ordering the layers topologically after removing feedback edges~\cite{eetemadi2019genetic}; knowledge-primed and pathway-informed networks predict higher-level phenotypes from curated structure~\cite{fortelny2020knowledge,lotfollahi2023biologically,elmarakeby2021biologically,selby2025beyond}. iPerturb keeps feedback and solves it to a unique convergence point, and it predicts genome-wide expression responses of unseen perturbations rather than a phenotype. Network-based simulation of perturbations has been used to follow transcription-factor knockouts through an inferred GRN by repeated matrix multiplication~\cite{kamimoto2023dissecting}; we add a sustained clamp, a signed steady state with a proven convergence point, and path-level evidence.

\paragraph{Learned perturbation-response models and their evaluation.} Latent-shift and compositional autoencoders~\cite{lotfollahi2019scgen,lotfollahi2023predicting}, graph models over gene--gene knowledge~\cite{roohani2024predicting}, and foundation or language models~\cite{cui2024scgpt,levine2024cell2sentence,adduri2025predicting} predict responses to unseen perturbations, but benchmarks show that they often do not surpass simple linear baselines~\cite{ahlmanneltze2025comparison} and that systematic shared variation can dominate standard metrics~\cite{vinas2025systema}. We therefore evaluate against a linear low-rank baseline, report centred and magnitude-residualised correlations, and use matched degree-preserving controls for every graph.

\paragraph{Network propagation.} Propagating labels or scores over biological networks, by random walks with restart, diffusion kernels or truncated walks, amplifies weak signals by integrating all paths between genes and down-weights paths through hubs~\cite{cowen2017network}. iPerturb uses propagation in two ways: signed, clamped propagation of the perturbation itself over the regulatory network (the explanation channel), and propagation of observed training responses over the association network (the predictive channel). Learnable univariate functions in a basis expansion, whose evaluation reduces to matrix multiplication, are used in learnable-activation networks~\cite{liu2024kan,li2024fastkan}; we use the same idea to keep Hill-shaped transfer functions monotone, bounded and cheap, in particular for edge-specific shapes.
